\documentclass[11pt]{article}
\usepackage{amsmath,amssymb,amsthm}
\usepackage[margin=1in]{geometry}
\newtheorem{proposition}{Proposition}
\newtheorem{corollary}{Corollary}
\newtheorem{definition}{Definition}

\title{History-Independent Complementarity Is Impossible\\
in Sequential Binary Choice}
\author{Draft notes}
\date{\today}

\begin{document}
\maketitle

\section{Model}

Two goods $A,B\in\{0,1\}$ have utilities $v_A$ and $v_B$. The parameter
$\gamma$ measures complementarity. Choosing $A$ changes the incremental
utility of $B$ by $\gamma$. When $\gamma>0$, taking $A$ makes $B$ more
attractive. When $\gamma<0$, taking $A$ makes $B$ less attractive. When
$\gamma=0$, the goods are independent.

A sequential binary model chooses one good at a time. Two assumptions govern
how utilities enter each step.

\begin{definition}[Reference independence]
\label{def:ref}
Each binary step depends only on its own incremental utility relative to not
taking the good.
\end{definition}

Let $f$ map an incremental utility into a choice probability. For every
real incremental utility $x$, the link satisfies $0<f(x)<1$. Reference
independence says that the choice probability at each step is $f$ of that
step's incremental utility alone.

\begin{definition}[History independence]
\label{def:hist}
Complementarity is \emph{history independent} if both choice orders use the
same link $f$ and the same $\gamma$. Which good is chosen first does not
change the form of complementarity.
\end{definition}

Under reference independence and history independence, choosing $A$ first
gives
\begin{equation}\label{eq:seq}
P(A=1)=f(v_A),\qquad P(B=1\mid A=a)=f(v_B+\gamma a).
\end{equation}
The first equality is reference independence at the first step. The
incremental utility of $A$ is just $v_A$, so $P(A=1)$ depends only on $v_A$.
The second equality is reference independence at the second step together
with complementarity. Conditional on $A=a$, the incremental utility of $B$ is
$v_B+\gamma a$, so $P(B=1\mid A=a)$ depends only on that quantity. History
independence requires the reverse order to take the same form with the same
$f$ and $\gamma$,
\begin{equation}\label{eq:rev}
P(B=1)=f(v_B),\qquad P(A=1\mid B=b)=f(v_A+\gamma b).
\end{equation}
From \eqref{eq:seq}, the joint probability of the bundle is therefore
\begin{equation}\label{eq:joint}
P(A=1,B=1)=f(v_A)\,f(v_B+\gamma).
\end{equation}

\section{The functional equation}

Bayes' identity holds for any joint distribution,
\[
P(A=1)\,P(B=1\mid A=1)=P(B=1)\,P(A=1\mid B=1).
\]
Substituting \eqref{eq:seq} and \eqref{eq:rev} forces
\begin{equation}\label{eq:FE}
f(v_A)\,f(v_B+\gamma)=f(v_B)\,f(v_A+\gamma)
\quad\Longleftrightarrow\quad
\frac{f(v_A+\gamma)}{f(v_A)}=\frac{f(v_B+\gamma)}{f(v_B)}.
\end{equation}
The multiplicative effect of $\gamma$ must not depend on the baseline
utility.

\section{Impossibility}

\begin{proposition}\label{prop:exp}
If $f>0$ satisfies \eqref{eq:FE} for all $v_A,v_B,\gamma$ and is continuous
(or measurable, monotone, or bounded on an interval), then $f(x)=Ce^{kx}$.
\end{proposition}

\begin{proof}
By \eqref{eq:FE} the ratio $f(x+\gamma)/f(x)$ is independent of $x$, so
$f(x+\gamma)=R(\gamma)f(x)$. With $g=\ln f$ and $h=\ln R$,
\[
g(x+\gamma)=g(x)+h(\gamma).
\]
Applying twice with increments $\gamma,\delta$ gives
$h(\gamma+\delta)=h(\gamma)+h(\delta)$, so $h$ is additive. Under regularity
$h(\gamma)=k\gamma$, hence $g(x)=g(0)+kx$ and $f(x)=Ce^{kx}$.
\end{proof}

Utilities are unbounded, so a probability link must satisfy $0<f(x)<1$ for
all $x$. The exponential $Ce^{kx}$ with $k\neq0$ is unbounded, so it fails.
The only bounded solution is $k=0$, so $f$ is constant. A constant link gives
$P(B=1\mid A=1)=P(B=1\mid A=0)$, so $\gamma$ has no effect.

\begin{corollary}\label{cor:imposs}
Under reference independence and history independence, no non-constant
bounded probability link $f$ can satisfy \eqref{eq:FE} with $\gamma\neq0$.
Either $\gamma=0$, or $f$ is constant and $\gamma$ has no effect on choice
probabilities.
\end{corollary}

The same impossibility has a direct reading through reference independence.
Complementarity requires the value of $A$ to depend on $B$, so the marginal
for $A$ must depend on $v_B$ and $\gamma$. Reference independence forbids
that dependence, because each step uses only its own incremental utility.
History independence and reference independence are therefore two faces of
the same constraint. The sequential decomposition forces each step to be a
standalone binary comparison, and that structure cannot carry
$\gamma\neq0$ in both orders.

\section{Why the joint bundle logit escapes}

The bundle model treats $\{0,A,B,AB\}$ as joint alternatives with utilities
\[
u_0=0,\quad u_A=v_A+\varepsilon_A,\quad u_B=v_B+\varepsilon_B,\quad
u_{AB}=v_A+v_B+\Gamma+\varepsilon_{AB},
\]
where the $\varepsilon_j$ are i.i.d.\ Type-I extreme value. Then
\[
P_{AB}=\frac{e^{v_A+v_B+\Gamma}}{1+e^{v_A}+e^{v_B}+e^{v_A+v_B+\Gamma}},
\]
and analogously for the other bundles. The derived conditional is
\[
P(B=1\mid A=a)=\Lambda(v_B+\Gamma a),
\]
but the derived marginal is not $\Lambda(v_A)$,
\[
P(A=1)=\frac{e^{v_A}\bigl(1+e^{v_B+\Gamma}\bigr)}
{1+e^{v_A}+e^{v_B}+e^{v_A+v_B+\Gamma}}.
\]

History independence \eqref{eq:FE} requires a single link $f$ to govern both
the marginal and the conditional in both orders. The bundle model never
imposes this. Its conditional is logistic, but its marginal is not
$\Lambda(v_A)$ for the same $\Lambda$. Equivalently, it violates reference
independence at the marginal, because $P(A=1)$ depends on $v_B$ and
$\Gamma$. That dependence is exactly what allows $\Gamma\neq0$. The
functional equation never arises, so no contradiction occurs.

\end{document}
